\section{Exploratory sensitivity to policy and investigation unit}
\label{sec:sensitivity}

\subsection{Mapping, shrinkage, severity, and reliability target}
The dated extension evaluates 23 configurations, including reproduction of the original setting, in 184 outer folds. The analysis retains all formula grids and operational references and repeats training-only mapping, cross-fitted reliability, and nested selection. It excludes the two learned families to focus on how explicit scoring formulas respond to input-policy changes. It does not assess the policy robustness of L1 or L2. Precision is varied over 0.3--0.7, purity over 0.5, 0.7, and 0.9, with four joint precision--purity corners. Support minima, shrinkage, severity, and criticality are varied separately (Online Resource~1, Section~S8). These settings are sensitivity cases, not a search for a new policy using test outcomes.

All nine class-level formulas still cover eight intervals in every configuration, although changing the criticality definition changes the denominator. Precision alone leaves S1\key{} at 21/32 intervals; raising purity to 0.9 reduces both S1\key{} and \Rsev{} to 8/32, including both joint corners with that purity. Shrinkage also matters: S1\key{} covers 23 intervals at $\lambda=0$ or 10 and 21 at 50 or 200, whereas \Rsev{} covers 11, 8, 18, and 21 respectively. A $+1$ shift to ranking severities with evaluation labels fixed reduces \Rsev{} to 8/32 while S1\key{} remains at 21/32. Thus the class-level coverage equality persists across these cases, but neither identifier-level gains nor their relative burden are independent of the author-defined policy.

The original reliability target is membership in any attack interval, whereas the primary endpoint concerns critical intervals. An ablation changes only the class/key reliability counts and their priors to critical membership, retaining the original phase mapping. P then defers 1,907 critical records beyond half the queue, compared with 9,047 originally, while S1 defers 9,092; both still cover eight intervals. S1\key{} covers 23 intervals but defers 8,966 records, versus 21 intervals and 1,802 deferred originally. Target alignment therefore changes deep-queue retrieval in a formula-dependent way and does not guarantee improvement. Precision and purity failures overlap; the original rejection counts alone cannot identify a single causal cutoff.

\subsection{Causal randomized ties}
For seven fixed original rankers, 100 seeded random priorities replace timestamp priority within equal-score backlog blocks. Priorities are deterministic functions of each alert's stable hash and seed, so they can be assigned on arrival without future information. Arrivals, scores, the 12/h service rate, and horizons are unchanged. Table~\ref{tab:review-random} summarizes totals across eight scenarios per seed. Seed means are first calculated within each scenario for exploratory paired tests; seeds do not enlarge the sample size.

S1\key{} averages 30.76 covered intervals and 282.34 critical starts; 96.93\% of critical records remain uninvestigated. P averages 13.98 covered intervals. The direction of their coverage--burden contrast persists (both scenario-paired exploratory $p=0.0078$), but severity-only coverage changes substantially, from 8 to a mean of 17.19. For severity-only ranking (O\textsubscript{sev}), timestamp-first ties yield coverage of 8, below every tested random-tie total both statically (10--14) and in replay (14--21). Comparisons against this reference therefore depend materially on tie policy; these observations do not establish a worst case over all possible tie orders. Random ties change comparisons with severity-only ranking, but the raw-record gap between coverage and retrieval remains. These results apply to this randomized tie policy, not to all causal schedulers.

\begin{tableorg}[htbp]
\centering\small
\caption{Raw-alert replay at 12/h: 11,508 starts per ranker and seed. FIFO ties uses timestamp-first ordering; remaining columns give random-tie means over 100 seed totals and coverage range (of 32). Before end counts first starts before the exclusive interval end. Online Resource~1, Section~S9 also reports noncritical interval-positive starts.}
\label{tab:review-random}
\resizebox{\textwidth}{!}{\begin{tabular}{lrrrrr}
\toprule
Ranker & FIFO ties & Random mean [range] & Before end & Critical starts & \shortstack{Interval-negative\\starts} \\
\midrule
P & 13 & 13.98 [13--14] & 6.83 & 64.6 & 9275.7 \\
S1 & 13 & 13.98 [13--14] & 6.83 & 64.6 & 9275.7 \\
S0\key & 28 & 28.63 [28--29] & 6.48 & 182.5 & 9335.7 \\
S1\key & 30 & 30.76 [30--31] & 7.57 & 282.3 & 10970.4 \\
S2\key & 29 & 29.63 [29--30] & 6.48 & 257.5 & 9985.8 \\
\Rsev & 25 & 25.63 [25--26] & 6.48 & 113.5 & 9298.8 \\
O\textsubscript{sev} & 8 & 17.19 [14--21] & 8.09 & 129.0 & 11250.7 \\
\bottomrule
\end{tabular}
}
\end{tableorg}

\subsection{Closed-group investigation}
Alerts are grouped by detector, native key, and host in windows anchored at the first member (60, 300, or 900 seconds). A group becomes selectable at window closure and takes the maximum member score; no member arrives at or after closure. The same analyst starts at most 12 group investigations per hour. The observation horizon remains the raw stream's first-to-last alert extent, without draining; groups closing later remain uninvestigated. Labels enter only after selection. Groups partition into critical (any critical member), noncritical-positive (a positive member but no critical member), and interval-negative (no positive member), including groups with mixed labels. An investigated group covers the union of its members' critical intervals at its service start, so future members cannot confer earlier coverage. Equal service per closed group is a workload assumption, not measured SOC productivity. Service time is held fixed across window lengths even though longer windows can place more records in each group, so the per-window comparison changes the workload represented by one investigation.

At 300 seconds, 2,655,821 records form 43,439 groups, of which 611 contain a critical record. S1\key{} covers 30 intervals, investigates 163/611 critical groups (26.68\%), and reaches 7 intervals before their end (Table~\ref{tab:review-grouped}). P covers 27 intervals, compared with 13 under raw-alert service. S0\key{} and \Rsev{} have identical displayed outcomes at 300 seconds, as verified in the saved per-scenario results; equality of these summaries does not establish identical selected groups. At 60/900 seconds, S1\key{} still covers 30 intervals and retrieves 14.49/37.05\% of critical groups. The investigation unit therefore changes both the between-ranker coverage gap and the retrieval denominator. The raw 97.2\% uninvestigated figure must not be read as a fraction of unresolved incidents. Grouping reduces that fraction under the stated model, while coverage and timely investigation remain distinct outcomes.

\begin{tableorg}[htbp]
\centering\small
\caption{Grouped replay with 300-second windows and 12 groups/h: 11,500 starts per ranker. Coverage is out of 32 intervals. Starts partition into interval-negative, noncritical-positive, and critical groups (Online Resource~1, Section~S9). Critical groups reports investigated / all critical groups, including those not released by the horizon. Critical members counts raw critical records contained in investigated groups; these are not separate investigations or confirmed attacks.}
\label{tab:review-grouped}
\resizebox{\textwidth}{!}{\begin{tabular}{lrrrrrr}
\toprule
Ranker & Covered & Before end & Critical groups & Retrieved (\%) & \shortstack{Interval-negative\\groups} & Critical members \\
\midrule
P & 27 & 5 & 113/611 & 18.49 & 11,007 & 259 \\
S1 & 27 & 5 & 113/611 & 18.49 & 11,071 & 259 \\
S0\key & 30 & 6 & 174/611 & 28.48 & 10,902 & 363 \\
S1\key & 30 & 7 & 163/611 & 26.68 & 11,250 & 260 \\
S2\key & 30 & 7 & 174/611 & 28.48 & 10,923 & 363 \\
\Rsev & 30 & 6 & 174/611 & 28.48 & 10,902 & 363 \\
O\textsubscript{sev} & 8 & 0 & 18/611 & 2.95 & 11,482 & 29 \\
\bottomrule
\end{tabular}
}
\end{tableorg}
